% Extracción quirúrgica del propietario V2 05b: líneas 100--359.
% El resto permanece en este archivo como procedencia, pero no se publica.
\iffalse
% Macroestructura solenoidal estrictamente interna del continuo discreto.
% Este módulo no introduce ninguna ecuación clásica ni un problema exterior.

\chapter{Macroestructura interna de balance, memoria y telescopía}
\label{chap:pdfv2-sol-interna}

\begin{statusbox}
La única entrada de este capítulo es la estructura discreta del continuo.
La rama \(\mathrm{sol}\) se construye antes de cualquier evaluación escalar
exterior y antes de cualquier formulación diferencial clásica.  Su cadena
completa es
\[
 \mathcal C_{\rm cont}^{\rm disc}
 \supset\mathcal C_{\rm cont}^{\rm sol}
 \xrightarrow{\operatorname{Res}_{\rm sol}}
 \mathcal G_{\rm sol}
 \xrightarrow{\operatorname{pr}_{X,{\rm sol}}}
 X_{\chi,{\rm sol}}
 \xrightarrow{\mathcal A_{\rm sol}}
 \mathfrak M_{\rm sol}
 \xrightarrow{\mathcal P_{\rm sol}}
 \mathcal C_{\rm sol}^{\rm HMT}.
\]
Cada flecha conserva como coordenadas la historia y el aparato que la
produce.  Ninguna imagen de esta cadena es un cociente que olvide su fuente.
\end{statusbox}

\section{Dominio propio antes de toda evaluación posterior}
\label{sec:pdfv2-sol-dominio}

Sea
\[
 x=(x_0\xrightarrow{e_0}x_1\xrightarrow{e_1}\cdots)
 \in\mathcal C_{\rm cont}^{\rm disc}
\]
una historia superviviente.  Cada arista \(e_j\) conserva sus dos hojas,
orientación, residuo, cociente y carry entero.  Sean
\(N_j^+(e_j)\) y \(N_j^-(e_j)\) las multiplicidades de las dos hojas y sea
\(\Delta_{\gamma_j}\in\mathbb Z\) el coeficiente orientado de la ruta, con el
carry de la composición ya incorporado.  El incremento firmado interno es
\begin{equation}\label{eq:pdfv2-sol-b}
 b_j=b(e_j):=
 \Delta_{\gamma_j}\bigl(N_j^+(e_j)-N_j^-(e_j)\bigr)\in\mathbb Z.
\end{equation}
No se obtiene evaluando una trayectoria exterior: es una coordenada del
ledger de la propia arista.

Sus partes positiva y negativa se construyen por
\begin{equation}\label{eq:pdfv2-sol-bpm}
 b_j^+:=\frac{|b_j|+b_j}{2},
 \qquad
 b_j^-:=\frac{|b_j|-b_j}{2}.
\end{equation}
Por tanto,
\begin{equation}\label{eq:pdfv2-sol-bpm-identidades}
 b_j=b_j^+-b_j^-,
 \qquad
 |b_j|=b_j^++b_j^-,
 \qquad
 b_j^+b_j^-=0.
\end{equation}

Las dos memorias primitivas y sus dos lectores son
\begin{align}
 M_0^+=M_0^-&=0,
 &M_j^+&=\sum_{0\leq i<j}b_i^+,
 &M_j^-&=\sum_{0\leq i<j}b_i^-;
 \label{eq:pdfv2-sol-memorias}\\
 D_j^0&:=M_j^+-M_j^-,
 &H_j^0&:=M_j^++M_j^-.
 \label{eq:pdfv2-sol-dh}
\end{align}
En consecuencia,
\begin{equation}\label{eq:pdfv2-sol-dh-diferencias}
 D_{j+1}^0-D_j^0=b_j,
 \qquad
 H_{j+1}^0-H_j^0=|b_j|.
\end{equation}
La memoria firmada \(D^0\) y la memoria total \(H^0\) permanecen como
objetos distintos; conservar solamente su diferencia final destruiría la
genealogía.

La subestructura \(\mathcal C_{\rm cont}^{\rm sol}\) consta de las historias
para las que, en todo prefijo y todo refinamiento nonádico:
\begin{enumerate}
 \item existen los incrementos \eqref{eq:pdfv2-sol-b} y las dos memorias
       \eqref{eq:pdfv2-sol-memorias};
 \item los cilindros finos se agrupan en fibras de nueve extensiones con la
       misma ley de probabilidad proyectiva;
 \item las proyecciones de profundidad y resolución definidas abajo son
       compatibles con las extensiones de historia;
 \item cada fila de pérdidas y cada columna terminal es finita a horizonte
       finito;
 \item todo prefijo admite al menos una extensión compatible a cualquier
       horizonte posterior.
\end{enumerate}
Estas son condiciones de pertenencia a la rama, no premisas silenciosas
atribuidas a una historia arbitraria.

\fi
\section{Refinamiento nonádico común: inclusión, esperanza y memoria de balance}
\label{sec:pdfv2-sol-e9j9}

Sea \(Y_j(x)\) el conjunto finito de cilindros compatibles con el prefijo
\(x_{\leq j}\), provisto de la medida proyectiva \(\mu_j\).  La truncación
\[
 \tau_{j+1,j}:Y_{j+1}(x)\longrightarrow Y_j(x)
\]
tiene nueve extensiones admisibles sobre cada cilindro.  Se consideran los
espacios finitos
\[
 \mathscr H_j(x):=\ell^2\bigl(Y_j(x),\mu_j\bigr).
\]
La inclusión y la esperanza nonádicas son
\begin{align}
 J_{9,j}:\mathscr H_j&\longrightarrow\mathscr H_{j+1},
 &(J_{9,j}f)(y)&=f(\tau_{j+1,j}y),
 \label{eq:pdfv2-sol-j9}\\
 E_{9,j}:\mathscr H_{j+1}&\longrightarrow\mathscr H_j,
 &(E_{9,j}g)(z)&=\frac1{9}
   \sum_{\tau_{j+1,j}y=z}g(y).
 \label{eq:pdfv2-sol-e9}
\end{align}
La compatibilidad proyectiva de \(\mu_{j+1}\) y \(\mu_j\) da
\begin{equation}\label{eq:pdfv2-sol-ej-projector}
 E_{9,j}J_{9,j}=I_{\mathscr H_j},
 \qquad
 J_{9,j}^*=E_{9,j},
 \qquad
 P_{9,j}:=J_{9,j}E_{9,j}=P_{9,j}^*=P_{9,j}^2.
\end{equation}
El proyector microscópico complementario es
\begin{equation}\label{eq:pdfv2-sol-qmicro-preliminar}
 Q_{9,j}:=I-P_{9,j}.
\end{equation}

Si \(b_f\) es el incremento sobre la fibra fina y
\begin{equation}\label{eq:pdfv2-sol-bcoarse}
 b_c:=E_{9,j}b_f,
\end{equation}
la convexidad elemental de \(|\cdot|\) construye, no postula, el carry de
cancelación
\begin{equation}\label{eq:pdfv2-sol-kappa}
 \kappa_j
 :=\frac{E_{9,j}|b_f|-|b_c|}{2}\geq0.
\end{equation}
En efecto, usando \(t^\pm=(|t|\pm t)/2\), se obtiene
\begin{align}
 E_{9,j}b_f^+
 &=\frac{E_{9,j}|b_f|+E_{9,j}b_f}{2}
   =b_c^++\kappa_j,
 \label{eq:pdfv2-sol-jensen-plus}\\
 E_{9,j}b_f^-
 &=\frac{E_{9,j}|b_f|-E_{9,j}b_f}{2}
   =b_c^-+\kappa_j.
 \label{eq:pdfv2-sol-jensen-minus}
\end{align}
Las dos hojas pierden la misma cantidad al comprimirse; por eso la diferencia
firmada se conserva mientras la memoria total registra la cancelación:
\begin{equation}\label{eq:pdfv2-sol-jensen-dh}
 E_{9,j}(b_f^+-b_f^-)=b_c,
 \qquad
 E_{9,j}(b_f^++b_f^-)=|b_c|+2\kappa_j.
\end{equation}

\section{Factorización ortogonal de los grados de profundidad y de las capas de resolución en \(\mathscr H_j^{\mathrm{depth}}\widehat\otimes\mathscr H_j^{\mathrm{res}}\)}
\label{sec:pdfv2-sol-micro-shell}

Para impedir que una misma pérdida sea contada a la vez como profundidad y
como resolución, el espacio de cada nivel se tipa como
\begin{equation}\label{eq:pdfv2-sol-factor-hilbert}
 \mathscr H_j
 =\mathscr H_j^{\rm depth}\widehat\otimes
  \mathscr H_j^{\rm res}.
\end{equation}
Para \(j\geq1\), el primer factor porta
\[
 P_j^{\rm depth}:=J_{9,j-1}E_{9,j-1}
 \quad\text{sobre }\mathscr H_j^{\rm depth},
\]
y se fija \(P_0^{\rm depth}=I\). El segundo factor lleva la partición
finita de shells heredada de la frontera de cilindros.  Si
\(\sigma_j:Y_j^{\rm res}\to S_j\) es esa partición, sus operadores son
\begin{align}
 J_j^{\rm res}f&=f\circ\sigma_j,
 &E_j^{\rm res}g(s)&=
 \frac1{\#\sigma_j^{-1}(s)}
 \sum_{\sigma_j(y)=s}g(y),
 \label{eq:pdfv2-sol-shell-je}\\
 P_j^{\rm res}&:=J_j^{\rm res}E_j^{\rm res}.
 \label{eq:pdfv2-sol-shell-p}
\end{align}

Se definen los tres proyectores mutuamente ortogonales
\begin{align}
 P_j^{\rm term}
 &:=P_j^{\rm depth}\otimes P_j^{\rm res},
 \label{eq:pdfv2-sol-pterm}\\
 Q_j^{\rm micro}
 &:=(I-P_j^{\rm depth})\otimes I,
 \label{eq:pdfv2-sol-qmicro}\\
 Q_j^{\rm shell}
 &:=P_j^{\rm depth}\otimes(I-P_j^{\rm res}).
 \label{eq:pdfv2-sol-qshell}
\end{align}
Entonces
\begin{equation}\label{eq:pdfv2-sol-ortogonalidad}
 Q_j^{\rm micro}Q_j^{\rm shell}=0,
 \qquad
 I=P_j^{\rm term}+Q_j^{\rm micro}+Q_j^{\rm shell},
\end{equation}
y, para todo \(\xi\in\mathscr H_j\),
\begin{equation}\label{eq:pdfv2-sol-pitagoras}
 \|\xi\|^2
 =\|P_j^{\rm term}\xi\|^2
  +\|Q_j^{\rm micro}\xi\|^2
  +\|Q_j^{\rm shell}\xi\|^2.
\end{equation}
No aparece el término
\((I-P_j^{\rm depth})\otimes(I-P_j^{\rm res})\) por segunda vez: ya está
contenido exactamente una vez en \(Q_j^{\rm micro}\).

\section{Extensiones de cilindro y fila completa de pérdidas}
\label{sec:pdfv2-sol-gamma-loss}

Una extensión de cilindros de nivel \(j\) a nivel \(j+1\) induce el mapa
\begin{equation}\label{eq:pdfv2-sol-gamma}
 \Gamma_j:\mathscr H_j\longrightarrow\mathscr H_{j+1},
 \qquad
 \Gamma_j f=f\circ\tau_{j+1,j},
\end{equation}
con todas las etiquetas de hoja, orientación, carry, memoria y ruta
transportadas en la base de cilindros.  Para \(n>j\), se escribe
\begin{equation}\label{eq:pdfv2-sol-gamma-composition}
 \Gamma_{j:n}:=\Gamma_{n-1}\cdots\Gamma_j,
 \qquad
 \Gamma_{j:j}:=I.
\end{equation}

Sobre \(\mathscr H_j\), los incrementos no negativos definen operadores de
multiplicación
\begin{equation}\label{eq:pdfv2-sol-loss-multipliers}
 B_j^+:=M_{\sqrt{b_j^+}},
 \qquad
 B_j^-:=M_{\sqrt{b_j^-}},
 \qquad
 K_j:=M_{\sqrt{2\kappa_j}}.
\end{equation}
Cada raíz se toma punto a punto sobre el conjunto finito de cilindros; si una
coordenada es nula, su operador es nulo.  El espacio de pérdidas es la suma
ortogonal externa
\[
 \mathscr E_j
 :=\mathscr H_j^{(+)}\oplus\mathscr H_j^{(-)}
   \oplus\mathscr H_j^{(\kappa)}
   \oplus\mathscr H_j^{({\rm micro})}
   \oplus\mathscr H_j^{({\rm shell})},
\]
y la fila completa es
\begin{equation}\label{eq:pdfv2-sol-loss-row}
 L_j:\mathscr H_j\longrightarrow\mathscr E_j,
 \qquad
 L_j\xi
 =\bigl(
 B_j^+\xi,B_j^-\xi,K_j\xi,
 Q_j^{\rm micro}\xi,Q_j^{\rm shell}\xi
 \bigr).
\end{equation}
Por tanto,
\begin{equation}\label{eq:pdfv2-sol-loss-norm}
 \|L_j\xi\|^2
 =\|B_j^+\xi\|^2+\|B_j^-\xi\|^2+\|K_j\xi\|^2
  +\|Q_j^{\rm micro}\xi\|^2
  +\|Q_j^{\rm shell}\xi\|^2.
\end{equation}
Las partes positiva, negativa, carry de cancelación, microestructura y shell
aparecen una vez cada una.  No se fusionan en una constante anónima ni se
repiten en dos sumandos.

\section{Terminal de historia completa}
\label{sec:pdfv2-sol-terminal}

Sea \(\operatorname{Hist}_{J}^{\rm sol}\) el conjunto de historias
supervivientes hasta el horizonte \(J\), cada una con su paquete de trece
coordenadas.  El espacio terminal se define como el espacio libre de esos
registros completos:
\begin{equation}\label{eq:pdfv2-sol-terminal-space}
 \mathscr T_J^{\rm sol}
 :=\ell^2\!\left(
 \left\{(h,D_{{\rm sol},J}(h)):
 h\in\operatorname{Hist}_{J}^{\rm sol}\right\}
 \right).
\end{equation}
La evaluación terminal
\begin{equation}\label{eq:pdfv2-sol-terminal-evaluation}
 E_J^{\rm term}:\mathscr H_J\longrightarrow\mathscr T_J^{\rm sol}
\end{equation}
envía cada vector básico de cilindro al registro completo de la historia que
lo genera y se extiende linealmente.  En particular, el terminal retiene
fibra, relaciones, números, orientación, carry, memoria, frontera, ruta,
multisección, supervivencia y lector.  No se sustituye por el último valor
visible de la ruta.

\section{Columna de observación, Gram y telescopía}
\label{sec:pdfv2-sol-gram}

Para \(0\leq j\leq J\), la columna futura completa es
\begin{equation}\label{eq:pdfv2-sol-g-column}
 \begin{aligned}
 G_{j;J}:\mathscr H_j&\longrightarrow
 \mathscr T_J^{\rm sol}\oplus
 \bigoplus_{n=j}^{J-1}\mathscr E_n,\\
 G_{j;J}\xi
 &:={}
 E_J^{\rm term}\Gamma_{j:J}\xi
 \ \oplus\ 
 \bigoplus_{n=j}^{J-1}L_n\Gamma_{j:n}\xi.
 \end{aligned}
\end{equation}
Para \(j=J\), la suma de pérdidas es vacía y
\(G_{J;J}=E_J^{\rm term}\).  El Gram terminal es
\begin{equation}\label{eq:pdfv2-sol-u-gram}
 U_{j;J}:=G_{j;J}^*G_{j;J}\succeq0.
\end{equation}

Reordenando la suma ortogonal del codominio se obtiene la identidad literal
\begin{equation}\label{eq:pdfv2-sol-g-recursion}
 G_{j;J}\xi
 \simeq
 \bigl(L_j\xi,\,G_{j+1;J}\Gamma_j\xi\bigr).
\end{equation}
En consecuencia,
\begin{equation}\label{eq:pdfv2-sol-stein}
 \boxed{U_{j;J}
 =L_j^*L_j+\Gamma_j^*U_{j+1;J}\Gamma_j}
\end{equation}
y
\begin{equation}\label{eq:pdfv2-sol-stein-difference}
 U_{j;J}-\Gamma_j^*U_{j+1;J}\Gamma_j
 =L_j^*L_j\succeq0.
\end{equation}
No hay una mutación de signo: el término de pérdida está en el lado positivo
de la identidad porque procede de una suma de cuadrados.

Iterando \eqref{eq:pdfv2-sol-stein} se obtiene la telescopía exacta
\begin{equation}\label{eq:pdfv2-sol-telescopia-operatorial}
 U_{j;J}
 =\Gamma_{j:J}^*(E_J^{\rm term})^*E_J^{\rm term}\Gamma_{j:J}
  +\sum_{n=j}^{J-1}
   \Gamma_{j:n}^*L_n^*L_n\Gamma_{j:n},
\end{equation}
o, de forma cuadrática,
\begin{equation}\label{eq:pdfv2-sol-telescopia-normas}
 \langle U_{j;J}\xi,\xi\rangle
 =\|E_J^{\rm term}\Gamma_{j:J}\xi\|^2
  +\sum_{n=j}^{J-1}\|L_n\Gamma_{j:n}\xi\|^2.
\end{equation}
El primer término conserva el futuro terminal completo; la suma conserva,
nivel a nivel, todas las pérdidas y todos los carries que condujeron hasta él.

\iffalse
\section{Las trece coordenadas de la restricción}
\label{sec:pdfv2-sol-trece}

Para un prefijo \(x_{\leq j}\), se define
\begin{equation}\label{eq:pdfv2-sol-d-package}
 \begin{aligned}
 D_{{\rm sol},j}(x):=\bigl(&
 \mathcal F_{{\rm sol},j},
 \operatorname{Ext}_{\Gamma,{\rm sol},j},
 \operatorname{Rel}_{{\rm sol},j},
 \mathcal N_{{\rm sol},j},
 \operatorname{or}_{{\rm sol},j},
 \operatorname{Carr}_{{\rm sol},j},
 \operatorname{Mem}_{{\rm sol},j},\\
 &\partial_{{\rm sol},j},
 \gamma_{{\rm sol},j},
 \operatorname{Msect}_{{\rm sol},j},
 \operatorname{Surv}_{{\rm sol},j},
 \operatorname{Term}_{{\rm sol},j},
 \mathcal L_{{\rm sol},j}\bigr).
 \end{aligned}
\end{equation}
Sus componentes son las siguientes.
\begin{enumerate}
 \item \(\mathcal F_{{\rm sol},j}\) contiene los cilindros \(Y_j\), su medida
       proyectiva, \(\mathscr H_j^{\rm depth}\),
       \(\mathscr H_j^{\rm res}\), sus productos tensoriales y los cinco
       espacios de pérdidas.
 \item \(\operatorname{Ext}_{\Gamma,{\rm sol},j}\) contiene
       \(J_{9,j},E_{9,j},\Gamma_j,\Gamma_{j:n}\), las inclusiones de shell y
       todos los mapas de truncación y refinamiento.
 \item \(\operatorname{Rel}_{{\rm sol},j}\) contiene
       \eqref{eq:pdfv2-sol-bpm-identidades},
       \eqref{eq:pdfv2-sol-ej-projector},
       \eqref{eq:pdfv2-sol-jensen-plus}--
       \eqref{eq:pdfv2-sol-jensen-minus},
       \eqref{eq:pdfv2-sol-ortogonalidad} y
       \eqref{eq:pdfv2-sol-stein}.
 \item \(\mathcal N_{{\rm sol},j}\) conserva
       \(j,J,9,b_j,b_j^+,b_j^-,\kappa_j\), multiplicidades de hoja, tamaños
       de fibras y dimensiones de shells.
 \item \(\operatorname{or}_{{\rm sol},j}\) conserva
       \(\Delta_{\gamma_j}\), el orden de ruta y la acción de inversión
       \(b_j\mapsto-b_j\), que intercambia \(b_j^+\) y \(b_j^-\) sin cambiar
       \(|b_j|\) ni \(\kappa_j\).
 \item \(\operatorname{Carr}_{{\rm sol},j}\) conserva el carry TPK entero,
       el carry de composición y el carry de cancelación \(\kappa_j\) como
       coordenadas diferentes.
 \item \(\operatorname{Mem}_{{\rm sol},j}\) es
       \((M_j^+,M_j^-,D_j^0,H_j^0)\) junto con el ledger ordenado de todos los
       incrementos anteriores.
 \item \(\partial_{{\rm sol},j}\) contiene la frontera orientada de cilindros,
       los incrementos \(b_j\), la fila \(L_j\) y las diferencias de Stein.
 \item \(\gamma_{{\rm sol},j}\) es la ruta completa
       \(e_0e_1\cdots e_{j-1}\), no sólo sus extremos.
 \item \(\operatorname{Msect}_{{\rm sol},j}\) contiene todas las extensiones
       nonádicas y de shell compatibles; no elige una prolongación.
 \item \(\operatorname{Surv}_{{\rm sol},j}\) registra los horizontes
       \(J\geq j\) para los que existen \(G_{j;J}\), \(U_{j;J}\) y
       terminales compatibles.
 \item \(\operatorname{Term}_{{\rm sol},j}\) contiene
       \(E_J^{\rm term}\Gamma_{j:J}\), el registro completo de historia y
       las trece coordenadas en cada horizonte.
 \item \(\mathcal L_{{\rm sol},j}\) publica internamente
       \((b^\pm,M^\pm,D^0,H^0,\kappa,Q^{\rm micro},Q^{\rm shell},
       L,G,U)\) y sus identidades.
\end{enumerate}

Para una extensión \(u:x_j\to x_{j'}\), el transporte coordinado
\(u_{{\rm sol},j',j}\) satisface
\begin{equation}\label{eq:pdfv2-sol-naturality-d}
 u_{{\rm sol},j',j}D_{{\rm sol},j'}(x_{j'})
 =D_{{\rm sol},j}(\tau_{j',j}x_{j'}).
\end{equation}
La igualdad incluye los cinco componentes de \(L\), el terminal completo y
la composición de \(\Gamma\); no es una igualdad sólo de cardinales.

\section{Grafo de restricción y objeto dependiente}
\label{sec:pdfv2-sol-graphs}

Se define
\begin{equation}\label{eq:pdfv2-sol-g-graph}
 \mathcal G_{\rm sol}:=\operatorname{Graph}(D_{\rm sol}),
 \qquad
 \operatorname{Res}_{\rm sol}(x):=(x,D_{\rm sol}x).
\end{equation}
La primera proyección es inversa sobre el grafo:
\begin{equation}\label{eq:pdfv2-sol-res-inverse}
 \pi_1\operatorname{Res}_{\rm sol}=I,
 \qquad
 \operatorname{Res}_{\rm sol}\pi_1=I_{\mathcal G_{\rm sol}}.
\end{equation}

Sea \(\chi_{\rm sol}(x)\) la multisección completa formada por todos los
cilindros, refinamientos, shells, filas de pérdidas y terminales compatibles
con \(x\).  Entonces
\begin{equation}\label{eq:pdfv2-sol-xchi}
 X_{\chi,{\rm sol}}
 :=\{(\chi_{\rm sol}(x),x,D_{\rm sol}x):
 x\in\mathcal C_{\rm cont}^{\rm sol}\},
\end{equation}
y
\begin{equation}\label{eq:pdfv2-sol-prx}
 \operatorname{pr}_{X,{\rm sol}}(x,D_{\rm sol}x)
 :=(\chi_{\rm sol}(x),x,D_{\rm sol}x).
\end{equation}
De nuevo, la proyección que conserva \((x,D_{\rm sol}x)\) es inversa sobre
esta imagen dependiente.

\section{Ensamblador y publicador internos}
\label{sec:pdfv2-sol-assembler-publisher}

El ambiente \(\mathscr M_{\rm sol}^{\rm amb}\) consta de sistemas
compatibles
\[
 \bigl(
 \{b,b^\pm,M^\pm,D^0,H^0,\kappa\},
 \{E_9,J_9,P^{\rm term},Q^{\rm micro},Q^{\rm shell}\},
 \{\Gamma,L,E^{\rm term},G,U\}
 \bigr)
\]
con todas las identidades anteriores.  El ensamblador bruto es
\begin{equation}\label{eq:pdfv2-sol-assembler}
 \begin{aligned}
 a_{\rm sol}:X_{\chi,{\rm sol}}&\longrightarrow
 \mathscr M_{\rm sol}^{\rm amb},\\
 a_{\rm sol}(\chi(x),x,Dx)&:=
 \bigl(
 \{b,b^\pm,M^\pm,D^0,H^0,\kappa\}_x,
 \{E_9,J_9,P,Q\}_x,
 \{\Gamma,L,E^{\rm term},G,U\}_x
 \bigr).
 \end{aligned}
\end{equation}
La macroestructura conserva su antecedente mediante
\begin{equation}\label{eq:pdfv2-sol-macro-graph}
 \mathfrak M_{\rm sol}:=\operatorname{Graph}(a_{\rm sol}),
 \qquad
 \mathcal A_{\rm sol}(z):=(z,a_{\rm sol}z).
\end{equation}

Sea \(\mathscr Y_{\rm sol}^{\rm int}\) el tipo de certificados formados por
las identidades
\begin{equation}\label{eq:pdfv2-sol-certificate-list}
 \begin{gathered}
 b=b^+-b^-,\qquad |b|=b^++b^-,\\
 E_9J_9=I,\qquad J_9E_9=P_9,\\
 E_9b_f^\pm=b_c^\pm+\kappa,\qquad\kappa\geq0,\\
 I=P^{\rm term}+Q^{\rm micro}+Q^{\rm shell},
 \qquad Q^{\rm micro}Q^{\rm shell}=0,\\
 U_{j;J}=L_j^*L_j+\Gamma_j^*U_{j+1;J}\Gamma_j,\\
 \langle U_{j;J}\xi,\xi\rangle
 =\|E_J^{\rm term}\Gamma_{j:J}\xi\|^2
  +\sum_{n=j}^{J-1}\|L_n\Gamma_{j:n}\xi\|^2,
 \end{gathered}
\end{equation}
junto con la naturalidad de todas las flechas.  El publicador bruto
\begin{equation}\label{eq:pdfv2-sol-publisher}
 p_{\rm sol}^{\rm raw}:\mathfrak M_{\rm sol}
 \longrightarrow\mathscr Y_{\rm sol}^{\rm int}
\end{equation}
publica ese certificado sin eliminar el macroobjeto.  Finalmente,
\begin{equation}\label{eq:pdfv2-sol-output-graph}
 \mathcal C_{\rm sol}^{\rm HMT}
 :=\operatorname{Graph}(p_{\rm sol}^{\rm raw}),
 \qquad
 \mathcal P_{\rm sol}(m):=(m,p_{\rm sol}^{\rm raw}m).
\end{equation}

\section{Teorema interno de generación solenoidal}
\label{sec:pdfv2-sol-theorem}

\begin{theorem}[Generación interna de la macroestructura de balance]
\label{thm:pdfv2-sol-internal-generation}
Para todo \(x\in\mathcal C_{\rm cont}^{\rm sol}\), la cadena
\[
 \mathcal C_{\rm cont}^{\rm sol}
 \xrightarrow{\operatorname{Res}_{\rm sol}}
 \mathcal G_{\rm sol}
 \xrightarrow{\operatorname{pr}_{X,{\rm sol}}}
 X_{\chi,{\rm sol}}
 \xrightarrow{\mathcal A_{\rm sol}}
 \mathfrak M_{\rm sol}
 \xrightarrow{\mathcal P_{\rm sol}}
 \mathcal C_{\rm sol}^{\rm HMT}
\]
construye de forma natural una macroestructura que contiene las dos memorias,
el carry de cancelación, la separación ortogonal micro/shell, la fila completa
de pérdidas, la columna futura, el terminal de historia completa, el Gram y la
telescopía.  Las trece coordenadas del continuo pueden recuperarse por
proyecciones sucesivas desde cualquier punto de la cadena.
\end{theorem}

\begin{proof}
Las identidades \eqref{eq:pdfv2-sol-bpm-identidades} y
\eqref{eq:pdfv2-sol-dh-diferencias} se obtienen directamente de la
descomposición positiva y negativa del incremento entero; por ello las
memorias no son datos añadidos.  Las fórmulas
\eqref{eq:pdfv2-sol-j9}--\eqref{eq:pdfv2-sol-e9}, junto con la medida
proyectiva, prueban \(E_9J_9=I\), \(J_9^*=E_9\) y que \(P_9\) es un
proyector ortogonal.

La desigualdad triangular da
\(E_9|b_f|\geq|E_9b_f|=|b_c|\), de modo que
\(\kappa\geq0\).  Sustituir
\(b^\pm=(|b|\pm b)/2\) prueba
\eqref{eq:pdfv2-sol-jensen-plus} y
\eqref{eq:pdfv2-sol-jensen-minus} sin una hipótesis adicional.

Las definiciones tensoriales
\eqref{eq:pdfv2-sol-pterm}--\eqref{eq:pdfv2-sol-qshell} prueban que
\(P^{\rm term}\), \(Q^{\rm micro}\) y \(Q^{\rm shell}\) son ortogonales y
suman la identidad.  En consecuencia, la fila \(L_j\) cuenta cada sector una
sola vez.

La definición \eqref{eq:pdfv2-sol-g-column} separa ortogonalmente el terminal
y las pérdidas de cada nivel.  Separar el primer sumando de pérdidas produce
\eqref{eq:pdfv2-sol-g-recursion}.  Tomar adjunto y componer prueba la identidad
de Stein \eqref{eq:pdfv2-sol-stein}; iterarla prueba
\eqref{eq:pdfv2-sol-telescopia-operatorial} y
\eqref{eq:pdfv2-sol-telescopia-normas}.  Todo término es una norma cuadrada,
por lo que no puede aparecer el signo opuesto.

La compatibilidad de truncación de cilindros, esperanza, proyecciones y
extensiones \(\Gamma\) da la naturalidad
\eqref{eq:pdfv2-sol-naturality-d}.  Finalmente,
\(\mathcal G_{\rm sol}\), \(\mathfrak M_{\rm sol}\) y
\(\mathcal C_{\rm sol}^{\rm HMT}\) son grafos, y
\(X_{\chi,{\rm sol}}\) es un objeto dependiente que conserva explícitamente
su antecedente.  La primera proyección de cada etapa recupera la etapa
anterior.  Así ninguna composición borra historia, memoria, multisección,
terminal ni lector.
\end{proof}

\section{Procedencia probatoria y falsadores}
\label{sec:pdfv2-sol-provenance-falsifiers}

\paragraph{Resultado recuperado.}
Pertenecen a la arquitectura ya construida el incremento firmado con carry,
la separación \(b=b^+-b^-\), las memorias \(M^\pm,D^0,H^0\), el par
\(E_9,J_9\), el carry de Jensen \(\kappa\), la extensión de cilindros, la
fila de pérdidas, el terminal, el Gram, Stein y la telescopía.

\paragraph{Formalización nueva.}
Son formalizaciones explícitas de esa arquitectura el tipado previo a toda
evaluación exterior, la factorización
\(\mathscr H^{\rm depth}\widehat\otimes\mathscr H^{\rm res}\), la separación
ortogonal que impide el doble conteo, el terminal libre sobre historias con
trece coordenadas y los cuatro empaquetados no ablativos por grafos.

\begin{criterion}[Falsadores internos de la rama \(\mathrm{sol}\)]
\label{crit:pdfv2-sol-falsifiers}
La construcción falla si ocurre cualquiera de los hechos siguientes:
\begin{enumerate}
 \item \(b\neq b^+-b^-\) o \(|b|\neq b^++b^-\);
 \item \(E_9J_9\neq I\), la medida no es proyectiva o \(P_9\) no es
       proyector;
 \item \(\kappa<0\) o fallan
       \(E_9b_f^\pm=b_c^\pm+\kappa\);
 \item \(Q^{\rm micro}Q^{\rm shell}\neq0\), o un mismo sector aparece en
       ambos sumandos de la fila \(L\);
 \item una de \(b^+,b^-,\kappa,Q^{\rm micro},Q^{\rm shell}\) se omite o se
       cuenta dos veces;
 \item el terminal conserva sólo el valor visible y no la historia con sus
       trece coordenadas;
 \item falla la identidad de Stein o la telescopía contiene un signo
       contrario a una suma de cuadrados;
 \item una extensión no conmuta con \(E_9,J_9,L,G,U\);
 \item se sustituye la multisección de continuaciones por una prolongación
       escogida retrospectivamente;
 \item se omite cualquiera de las trece coordenadas.
\end{enumerate}
En el último caso el diagnóstico vinculante es: \emph{objeto sustituido;
conclusión no transferible}.
\end{criterion}

\begin{warningbox}
Este capítulo concluye únicamente la construcción interna de
\(\mathcal C_{\rm sol}^{\rm HMT}\).  No contiene una ecuación diferencial
clásica, datos iniciales clásicos, un mapa desde datos exteriores ni una
conclusión sobre un problema exterior.  Esos entrelazadores pertenecen a otro
documento y no gobiernan la existencia de esta macroestructura.
\end{warningbox}
\fi
